% 可数（此处有限）状态链的沟通类。箭头表示一步正概率转移（数值省略）。
% {1,2}：可离开、不封闭，状态暂留；{3,4}：封闭类，非周期（3有自环）；
% {5,6}：封闭类，5->6->5确定交替，周期2。手写TikZ源。
\begin{tikzpicture}[font=\figurefont\fontsize{8}{10}\selectfont,text=FigureInk,>=Stealth,
  st/.style={circle,draw=FigureStroke,fill=white,line width=.75pt,minimum size=7mm,inner sep=0pt},
  ed/.style={->,draw=FigureStroke,line width=.75pt},
  cls/.style={rounded corners=1.8pt,line width=.75pt},
  note/.style={font=\FigureNoteFont\fontsize{7.5}{9}\selectfont,text=FigureMuted,align=center}]
\fill[cls,FigureYellowFill] (-6mm,-8mm) rectangle (26mm,8mm);
\draw[cls,FigureYellow] (-6mm,-8mm) rectangle (26mm,8mm);
\fill[cls,FigureBlueFill] (42mm,6mm) rectangle (72mm,22mm);
\draw[cls,FigureBlue] (42mm,6mm) rectangle (72mm,22mm);
\fill[cls,FigureGreenFill] (42mm,-22mm) rectangle (72mm,-6mm);
\draw[cls,FigureGreen] (42mm,-22mm) rectangle (72mm,-6mm);
\node[st] (1) at (0,0) {$1$};
\node[st] (2) at (20mm,0) {$2$};
\node[st] (3) at (48mm,14mm) {$3$};
\node[st] (4) at (66mm,14mm) {$4$};
\node[st] (5) at (48mm,-14mm) {$5$};
\node[st] (6) at (66mm,-14mm) {$6$};
\draw[ed] (1) to[bend left=25] (2);
\draw[ed] (2) to[bend left=25] (1);
\draw[ed] (2) -- (3);
\draw[ed] (2) -- (5);
\draw[ed] (3) to[bend left=25] (4);
\draw[ed] (4) to[bend left=25] (3);
\draw[ed] (3) to[out=120,in=60,looseness=6] (3);
\draw[ed] (5) to[bend left=25] (6);
\draw[ed] (6) to[bend left=25] (5);
\node[note,anchor=north] at (10mm,-9mm) {不封闭的类\\状态暂留};
\node[note,anchor=west] at (73mm,14mm) {封闭类\\非周期};
\node[note,anchor=west] at (73mm,-14mm) {封闭类\\周期$2$};
\end{tikzpicture}
